%*******************************************************************************
%*********************************** First Chapter *****************************
%*******************************************************************************

\chapter{Introduction}  %Title of the First Chapter

\ifpdf
    \graphicspath{{Chapter1/Figs/Raster/}{Chapter1/Figs/PDF/}{Chapter1/Figs/}}
\else
    \graphicspath{{Chapter1/Figs/Vector/}{Chapter1/Figs/}}
\fi

\section{Problem Background}

The University of San Carlos (USC) Clinic plays a crucial role in ensuring the health and well-being of its students, faculty, and staff. However, the clinic's current Health Information System (HIS) is predominantly paper-based and fragmented, creating several challenges that impede the delivery of efficient, comprehensive, and high-quality healthcare services.

\subsection{Incomplete Student Health Records}
The clinic's existing system struggles to capture complete medical records, particularly for incoming first-year students. This gap in health histories can result in missed opportunities for preventive care, delays in diagnosis and treatment, and difficulties in tracking health trends across the student population. Comprehensive health records are essential for effective healthcare delivery, enabling the clinic to provide timely and appropriate interventions.

\subsection{Inefficient Consultation Management}
Managing a high volume of consultations, especially for minor ailments such as coughs and colds, is a significant challenge under the current system. This inefficiency leads to longer wait times, overcrowding, and frustration for both students and staff. Additionally, it delays the attention needed for more urgent medical issues, compromising the clinic's overall effectiveness.

\subsection{Medication Inventory Control Issues}
The existing manual inventory management system, while adhering to the First-In, First-Out (FIFO) principle, is prone to errors. These errors can lead to stockouts and the expiration of unused medications, impacting the timely provision of care and resulting in unnecessary costs due to wasted resources. A robust inventory control system is vital for maintaining adequate medication supplies and ensuring patient safety.

\subsection{Limited Reporting and Data Analysis}
The clinic relies on Google Sheets for daily treatment records, which, while convenient, limits data accuracy, reporting capabilities, and long-term data analysis. This limitation hampers the clinic's ability to generate comprehensive reports and gain valuable insights into clinic utilization and student health outcomes. Effective data analysis is crucial for informed decision-making and continuous improvement of healthcare services.

\subsection{Cumbersome Administrative Processes}
The manual process for issuing medical certificates is time-consuming and inefficient, creating delays for both clinic staff and students. This inefficiency can hinder students' ability to promptly return to their academic and extracurricular activities. Streamlined administrative processes are necessary to reduce staff workload and improve the overall patient experience.

\subsection{Inadequate Information Dissemination}
The clinic lacks a robust mechanism for disseminating health information and educational materials to the USC community across both the Downtown and Talamban campuses. This gap limits opportunities for preventive care, health promotion, and timely updates on relevant health issues. Effective communication channels are essential for fostering a well-informed and health-conscious university community.

\subsection{Ineffective Patient Feedback Collection}
The current method of collecting student feedback is cumbersome and does not lend itself to systematic analysis. This inefficiency makes it difficult to gauge student satisfaction, identify areas for improvement, and make data-driven decisions to enhance the quality of care provided by the clinic. A streamlined feedback collection process is critical for ensuring continuous improvement and high standards of patient care.

\section{Statement of the Problem}

The USC Clinic's reliance on a predominantly paper-based and fragmented patient information management system has resulted in several inefficiencies, including:
\begin{enumerate}
    \item Incomplete patient records
    \item Limited reporting capabilities
    \item Lack of streamlined communication channels for health information dissemination
    \item Inefficient patient feedback collection
\end{enumerate}

These challenges hinder the clinic's ability to provide timely, efficient, and comprehensive healthcare services to the USC community. Addressing these issues is critical to improving clinic efficiency, enhancing patient satisfaction, and ensuring the highest quality of care for all USC students, faculty, and staff.

The project aims to tackle these specific problems by implementing a modern, web-based PIS that focuses on capturing comprehensive health and dental records, facilitating effective health information dissemination, and collecting and analyzing patient feedback. This initiative will lay the groundwork for future enhancements and broader deployment, ultimately transforming the USC Clinic’s healthcare delivery.

\section{Goals and Objectives}

This project aims to design, develop, and implement a modern PIS for the USC Clinic to address the following specific objectives:
\begin{enumerate}
    \item \textbf{Capture comprehensive medical and dental records for students:} Streamline data collection to ensure complete and up-to-date medical and dental records for all students, including medical histories, treatment plans, and dental records.
    \item \textbf{Increase the reach of health awareness campaigns:} Use the PIS to promote health campaigns with dedicated pages featuring PubMats, banners, and other materials, encouraging proactive health management.
    \item \textbf{Facilitate student feedback collection and analysis:} Integrate digital evaluation forms to systematically collect and analyze patient feedback, enabling continuous improvement in healthcare services.
    \item \textbf{Generate Detailed Reports for Data Analysis and Decision-Making:} Implement a module for generating customizable reports on patient visits, treatment outcomes, and feedback analysis to aid decision-making.
    \item \textbf{Issue Notifications for Health Campaigns, and Clinic Updates:} Include a notification system for sending alerts on health campaigns and important clinic updates via email and in-app notifications.
    \item \textbf{Optimize the Issuance of Medical Certificates:} Streamline the process of issuing medical certificates by allowing quick access and input of relevant patient information, ensuring accuracy and efficiency.
\end{enumerate}

By achieving these objectives, this project aims to transform the USC Clinic's patient information management, leading to improved healthcare delivery, enhanced patient satisfaction, and increased operational efficiency for the clinic staff.

\section{Significance of the Study}

The successful implementation of a modernized PIS at the USC Clinic has the potential to transform healthcare delivery for the USC community. This project will directly benefit students by:
\begin{enumerate}
    \item \textbf{Improving access to care:} Streamlined data collection and notifications will reduce wait times and make it easier for students to manage their health.
    \item \textbf{Enhancing the student experience:} Faster service, personalized care, and access to their health information electronically will improve the overall student healthcare experience.
    \item \textbf{Promoting student well-being:} Access to educational materials and targeted health campaigns through the PIS will empower students to take a more active role in their health, leading to better health outcomes.
\end{enumerate}

The benefits extend to the clinic as well:
\begin{enumerate}
    \item \textbf{Increasing efficiency:} By optimizing processes and centralizing data, the new PIS will free up staff time, allowing them to focus on direct patient care.
    \item \textbf{Improving decision-making:} Accurate data and detailed reports will empower clinic leadership to make informed decisions regarding resource allocation and service improvement.
    \item \textbf{Reducing costs:} Streamlining workflows, reducing paperwork, and optimizing medication inventory management will contribute to cost savings.
\end{enumerate}

Beyond the immediate impact on the USC Clinic, this project can serve as a model for other university health centers seeking to modernize their systems and enhance student care. The lessons learned and best practices identified can be shared with the broader healthcare community, potentially contributing to improvements in patient care across diverse settings. 

Additionally, the data collected through the new PIS can be used for research, potentially leading to new insights into student health trends and contributing to the development of more effective healthcare interventions in academic settings.

\section{Scope and Limitations}

\subsection{Scope}
The proposed PIS for the USC Clinic is designed to modernize and enhance the clinic's healthcare delivery through the following core functionalities:
\begin{enumerate}
    \item \textbf{Comprehensive Health and Dental Records:}
    \begin{enumerate}
        \item Centralized and secure storage of all student health information, including medical histories, treatment plans, and dental records.
        \item Focus on capturing comprehensive records from the start of student enrollment, ensuring complete and up-to-date information.
        \item The system will focus on students, but is not limited to students since faculty and university staff will be filling up the same forms.
    \end{enumerate}
    \item \textbf{Health Information Dissemination:}
    \begin{enumerate}
        \item Utilization of targeted communication tools to disseminate health information and promote preventive care.
        \item Dedicated web pages within the system for health awareness campaigns, featuring PubMats, banners, and other educational materials.
    \end{enumerate}
    \item \textbf{Patient Feedback System:}
    \begin{enumerate}
        \item Integration of user-friendly digital evaluation forms to collect and analyze patient feedback.
        \item Automated reminders to encourage continuous patient input, enabling the clinic to identify areas for improvement.
    \end{enumerate}
    \item \textbf{Report Generation:}
    \begin{enumerate}
        \item The system will include a report generation module that allows clinic staff to create detailed, customizable reports on various metrics such as patient visits, treatment outcomes, and feedback analysis. These reports will be exportable to formats like PDF and Excel for further analysis and record-keeping.
    \end{enumerate}
    \item \textbf{Notifications:}
    \begin{enumerate}
        \item The system will feature a notification system that sends alerts and reminders for health campaigns, and important clinic updates via email and in-app notifications.
    \end{enumerate}
    \item \textbf{Medical Certificate Issuance Optimization:}
    \begin{enumerate}
        \item The PIS will optimize the process of issuing medical certificates by allowing clinic staff to quickly access and input relevant patient information and include all necessary details. This will reduce administrative workload and speed up the issuance process.
    \end{enumerate}
    \item \textbf{Platform:}
    \begin{enumerate}
        \item The system will be web-based, ensuring accessibility through standard web browsers.
    \end{enumerate}
    \item \textbf{Scalability:}
    \begin{enumerate}
        \item The system will be designed to scale for future iterations and broader deployment across other clinics as needed.
    \end{enumerate}
\end{enumerate}

\subsection{Limitations}
\begin{enumerate}
    \item \textbf{Scope of Records:} The initial focus will be on capturing comprehensive medical and dental records for students, aligning with the clinic's objective to gather extensive health information during enrollment phases.
    \item \textbf{User Adoption:} The success of the PIS implementation depends heavily on comprehensive training and ongoing support for clinic staff. Effective user adoption strategies are critical to ensure the system is used both reliably and effectively.
    \item \textbf{Resource Constraints:} As a student-led project with limited resources, the development team may need to prioritize certain features over others. Some functionalities, such as medication tracking, inventory management, may be deferred to future iterations.
    \item \textbf{Security Measures:} Initially, only basic security measures will be implemented to protect patient data, with plans to enhance security features in future updates. The system's user authentication will be basic and will not be able to verify whether a student is currently enrolled or has graduated.
    \item \textbf{Standalone System:} The system will be operated independently and will not be integrated with existing systems.
\end{enumerate}

\subsection{Collaboration and Alignment}

The project team will collaborate closely with the USC Clinic to ensure the system aligns with the clinic’s operational needs and objectives. This collaboration includes regular feedback sessions and iterative development to refine and enhance the system based on real-world use and requirements.

By addressing these scopes and limitations, the project team aims to deliver a functional, user-friendly, and scalable PIS that significantly improves the USC Clinic’s healthcare delivery, while also setting the stage for future enhancements and broader deployment.
